\section{Implementation Details}
\label{app:implementation-details}

\paragraph{Random seeds.}
For the main-result tables, we run each setting with five seeds: 0, 42, 1234, 3407, and 2024. Unless stated otherwise, auxiliary analyses outside the main-result tables use seed 42.

\paragraph{Hardware and datasets.}
All experiments use a GPU cluster with four NVIDIA RTX A5000 GPUs. For image classification, we follow \safefix~\citep{xu2025safefix}: a ResNet-18 classifier~\citep{he2016deep} on an ImageNet subset~\citep{deng2009imagenet} with 10 classes, 1,300 images per class, and an 8:1:1 train/validation/test split, denoted ImageNet10. The ImageNet10 classes are \texttt{backpack}, \texttt{barber chair}, \texttt{coffee mug}, \texttt{desk}, \texttt{electric guitar}, \texttt{park bench}, \texttt{pitcher}, \texttt{purse}, \texttt{rocking chair}, and \texttt{water bottle}. We also evaluate CelebA~\citep{liu2015deep} lipstick-wearing classification with a ResNet-18-based binary classifier and 80,000/10,000/10,000 train/validation/test images. For pose estimation, we follow HiBug2~\citep{chen2025hibug2} on COCO person keypoints~\citep{lin2014microsoft}, using Keypoint R-CNN~\citep{mask-rcnn-2017} with a ResNet-50-FPN backbone~\citep{he2016deep}, 20,000 source training images, and 1,173-image held-out validation and test splits. For object detection, we use KITTI~\citep{geiger2012we} Car and Pedestrian instances converted to YOLO format with 5,000/1,240/1,241 train/validation/test images, and train YOLOv8n.

\paragraph{Training and repair schedule.}
For ImageNet10 classification, the diagnosis-stage ResNet-18 model is trained with cross-entropy loss and Adam~\citep{kingma2015adam}, using learning rate $1\times10^{-4}$, weight decay $1\times10^{-3}$, batch size 64, and 30 epochs. During repair, we train with the selected source subset and the generated repair images for 15 epochs using learning rate $1\times10^{-5}$. For CelebA, COCO, and KITTI, the repair-stage epoch count is set to half of the corresponding diagnosis-stage training schedule, with the same repair-stage learning rate $1\times10^{-5}$.

\paragraph{Image generation.}
For the generated repair images used in the main results, we use \controlnet~\citep{zhang2023adding}, a conditional diffusion model based on Stable Diffusion 1.5~\citep{rombach2022high}. The model is conditioned on soft HED boundaries~\citep{xie2015holistically} extracted from source images and uses 30 DDIM inference steps~\citep{ddim}. The soft boundary condition helps preserve object structure during edits, which is useful for attribute changes such as recoloring, material changes, and weather or style edits.

\paragraph{Attribute proposal and assignment.}
We use vision-language models (VLMs) in two separate roles. First, a VLM proposes the semantic attribute vocabulary $\attrs$ for a task. For example, in KITTI object detection, GPT-5~\citep{openai2025gpt5} is used to define candidate visual attributes such as color, weather, scale, viewpoint, occlusion, truncation, and motion state. Second, a VLM assigns attribute values to source and generated images. In the experiments reported in this paper, we use Qwen3-VL~\citep{bai2025qwen3vl} for this assignment step. For the KITTI object-detection repair, color and weather are the attribute fields tied to the discovered repair slice.

\paragraph{Generation and assignment cost.}
Generating 1,000 images with ControlNet takes slightly over one hour on our cluster. Assigning attributes to these images with Qwen3-VL takes about 10--20 minutes, depending on the queried attribute set.

